\documentclass[11pt]{article}

\usepackage[T1]{fontenc}
\usepackage{amsmath,amssymb}
\usepackage{booktabs}
\usepackage[margin=1in]{geometry}
\usepackage[hidelinks]{hyperref}
\usepackage{listings}
\usepackage{tabularx}
\lstset{basicstyle=\ttfamily\small,breaklines=true,columns=fullflexible,
  literate={≤}{{$\le$}}1 {∅}{{$\varnothing$}}1}

\title{Proof Dossier: Certified Positive Separation}
\author{Project Koios Formal}
\date{}

\begin{document}
\maketitle

\section{Identity and status}

\begin{tabularx}{\textwidth}{@{}l>{\raggedright\arraybackslash}X@{}}
\toprule
Statement ID & \texttt{BG-THM-003} \\
Status & \texttt{proved} \\
Lean declaration & \path{ProjectKoios.BandGauge.Geometry.disjoint_of_pos_le_admissibleSetEDist} \\
Owning module & \path{ProjectKoios/BandGauge/Geometry/AdmissibleSets.lean} \\
\bottomrule
\end{tabularx}

\section{Exact mathematical statement}

Let $X$ be a metric space, let $A,B\subseteq X$, and let
$\delta_{\mathrm e}(A,B)$ be their extended admissible-set distance. If an extended
nonnegative number $\ell$ satisfies
\[
 0<\ell\le\delta_{\mathrm e}(A,B),
\]
then $A$ and $B$ are disjoint.

In the scientific application, the premise is a certified global lower bound on the
distance between the spectral and aligned-operator admissible sets in one candidate
parameter domain.

\section{Source provenance}

\begin{tabularx}{\textwidth}{@{}l>{\raggedright\arraybackslash}X@{}}
\toprule
Owner repository & \texttt{ksdft2effmass} \\
Navigation branch & \path{work/paper-3-higher-rank-extension} \\
Immutable commit & \path{99a46cb59d0c6f27ada7d709f194996ef9619d9f} \\
Path & \path{docs/publications/papers/Ragasa2026c/v1/manuscript.md} \\
Locator & definition lines 117--130; proposition lines 132--138; proof lines 140--145 \\
Dirty-state limitation & later owner worktree changes are excluded \\
\bottomrule
\end{tabularx}

The declaration \texttt{formalizes} the source implication from a certified positive
lower bound to disjointness.

\section{Assumptions and representation delta}

\begin{tabularx}{\textwidth}{@{}l>{\raggedright\arraybackslash}X>{\raggedright\arraybackslash}X@{}}
\toprule
Item & Lean encoding & Boundary \\
\midrule
Ambient domain & \texttt{[MetricSpace X]} & Matches the source metric assumption. \\
Set distance & \texttt{ENNReal} & Extends the source's real-valued nonempty distance. \\
Certificate & $0<\ell\le\delta_{\mathrm e}(A,B)$ & The theorem checks the implication, not the external computation of $\ell$. \\
Conclusion & \texttt{Disjoint A B} & Scoped to the declared sets and domain. \\
\bottomrule
\end{tabularx}

\section{Prior art and contribution boundary}

The implication is a classical consequence of the definition of set distance. Validated
global optimization and interval analysis supply \texttt{background-only} methods for
constructing rigorous lower bounds; accepted References-owned IDs remain unavailable.
The theorem is not novel. The project contribution is the explicit use of this standard
implication in the joint-adequacy decision workflow.

\section{Proof schematic}

\begin{center}
\begin{tabular}{c}
Assume $0<\ell\le\delta_{\mathrm e}(A,B)$ \\
$\Downarrow$ \\
Assume for contradiction that $x\in A\cap B$ \\
$\Downarrow$ \\
BG-THM-002 gives $\delta_{\mathrm e}(A,B)=0$ \\
$\Downarrow$ \\
$0<\ell\le0$, contradiction \\
$\Downarrow$ \\
$A\cap B=\varnothing$
\end{tabular}
\end{center}

\section{Human-readable proof}

Assume that a point $x$ belongs to both sets. BG-THM-002 then gives
$\delta_{\mathrm e}(A,B)=0$. The certified lower bound becomes $\ell\le0$, contradicting
$0<\ell$. Therefore no point belongs to both sets, so the sets are disjoint.

\section{Lean realization}

\begin{lstlisting}
theorem disjoint_of_pos_le_admissibleSetEDist {A B : Set X}
    {lower : ENNReal} (hpos : 0 < lower)
    (hlower : lower ≤ admissibleSetEDist A B) :
    Disjoint A B
\end{lstlisting}

The proof invokes BG-THM-002 and the set-theoretic characterization of disjointness.

\section{Trusted base and axiom report}

Lean \texttt{v4.35.0} and Mathlib revision
\path{56a005855ea17e0b0958f73ec801ae140b0a1679} report:

\begin{lstlisting}
'ProjectKoios.BandGauge.Geometry.disjoint_of_pos_le_admissibleSetEDist'
depends on axioms: [propext, Classical.choice, Quot.sound]
\end{lstlisting}

There is no \texttt{sorry}, \texttt{admit}, or repository-owned scientific axiom.

\section{AI and review record}

The statement design, Lean proof, and dossier prose were AI-generated or AI-assisted.
Lean checked the declaration and proof in the pinned environment. Repository checks pass.
No independent expert mathematical or physical review is recorded.

\section{Scientific and evidence boundaries}

The theorem does not compute $\ell$, validate an optimizer, or prove that a numerical
bound encloses the global infimum. It proves the exact consequence of a supplied bound.
Disjointness applies only to the declared candidate family, parameter domain, losses,
alignment contract, and thresholds. It does not prove physical inequivalence or exclude a
larger model family.

\section{Unresolved work}

Accepted canonical reference IDs and a checked certificate interface for externally
computed lower bounds remain unavailable.

\section{Review record}

On 18 October 2026, the Lean build, axiom print, and repository checks passed. Statement
status: \texttt{proved}. The dossier remains a working draft.

\end{document}
